\documentclass[12pt,a4paper]{article}

% Encoding and language
\usepackage[T1]{fontenc}
\usepackage[utf8]{inputenc}
\usepackage[english]{babel}

% Mathematics
\usepackage{amsmath,amssymb,amsthm}
\usepackage{mathtools}

% Tables and graphics
\usepackage{array,booktabs,multirow}
\usepackage{graphicx}

% Page layout
\usepackage{geometry}
\geometry{left=2.5cm,right=2.5cm,top=2.5cm,bottom=2.5cm}

% Hyperlinks
\usepackage{hyperref}
\hypersetup{
    colorlinks=true,
    linkcolor=blue,
    citecolor=blue,
    urlcolor=blue
}

% Theorems and definitions
\newtheorem{definition}{Definition}[section]
\newtheorem{theorem}{Theorem}[section]
\newtheorem{proposition}{Proposition}[section]
\newtheorem{remark}{Remark}[section]

% Operators
\DeclareMathOperator{\dist}{dist}
\DeclareMathOperator{\ReLU}{Re}
\newcommand{\R}{\mathbb{R}}
\newcommand{\N}{\mathbb{N}}
\newcommand{\G}{\mathbf{G}}
\newcommand{\A}{\mathcal{A}}
\newcommand{\D}{\mathcal{D}}
\newcommand{\U}{\mathcal{U}}
\newcommand{\M}{\mathcal{M}}
\newcommand{\T}{\mathcal{T}}
\newcommand{\I}{\mathcal{I}}
\newcommand{\Pset}{\mathcal{P}}

\title{All Types of GRA-Zeroing: Static, Cyclic, Chaotic and Hybrid \\[0.5em]
\large Set-Theoretic and Operator Extension}

\author{A.~I.~Ivanov\thanks{Institute of Dynamical Systems, Moscow, Russia; ivanov@ids.ru} \and 
P.~S.~Petrov\thanks{Laboratory of Complex Systems, St. Petersburg, Russia; petrov@lcss.spb.ru}}
\date{\today}

\begin{document}

\maketitle

\begin{abstract}
GRA-zeroing (Geometric Recursive Analytics) is a family of methods for minimizing deviations from goals (``foam'') in multilevel dynamical systems. This paper provides a complete description of the three fundamental GRA types (static, cyclic, chaotic) and their hybrid combinations, together with mathematical applicability criteria. The concept is then extended to a set-theoretic and operator formalism: GRA types are regarded as elements of an abstract set $\G = \{S, C, H\}$ on which Boolean operations (union, intersection, difference, symmetric difference, complement) are defined. Hybrid architectures are described by Cartesian products $\G^L$. The notion of a GRA-zeroing operator acting on the space of dynamical systems is introduced, and an algebra of such operators is constructed. Practical applications of the set-theoretic approach are discussed: automatic system classification, design of hierarchical controllers, verification, and adaptive regime switching. The material concludes with a final formula for selecting the GRA type under multiple constraints.

\textbf{Keywords:} GRA-zeroing, dynamical systems, attractors, foam, static GRA, cyclic GRA, chaotic GRA, hybrid GRA, set theory, operator analysis, hierarchical systems, AGI/ASI.
\end{abstract}

\tableofcontents
\newpage

\section{Introduction}

GRA-zeroing (Geometric Recursive Analytics) is a family of methods for minimizing ``foam'' $\Phi(x)$ (deviations from goals) in multilevel systems. Depending on the dynamic behavior of the system, three fundamental attractor types are distinguished, corresponding to three kinds of GRA:

\begin{itemize}
    \item \textbf{Static GRA} – target attractor: fixed point ($\Phi \to 0$).
    \item \textbf{Cyclic GRA} – target attractor: limit cycle ($\Phi(t)$ minimal but not zero, periodic).
    \item \textbf{Chaotic GRA} – target attractor: strange attractor (statistical state, minimization of mean fluctuations).
\end{itemize}

In addition, hybrid variants are possible (e.g., different hierarchy levels use different types).

In this article we not only describe each kind in detail with mathematical formulas and applicability criteria (Sections~\ref{sec:static}--\ref{sec:hybrid}), but also extend the concept to a set-theoretic and operator formalism (Sections~\ref{sec:set}--\ref{sec:applications}), which allows a uniform combination of strategies and the synthesis of complex controllers for artificial intelligence systems.

\section{Static GRA (Fixed Point)}
\label{sec:static}

\subsection{Mathematical Model}

The system is described by an ordinary differential equation:
\begin{equation}
    \dot{x} = f(x), \quad x \in \R^n.
\end{equation}
Let $x^*$ be a stationary point ($f(x^*) = 0$). The foam $\Phi(x)$ is a positive definite function, e.g., the squared distance to the goal:
\begin{equation}
    \Phi(x) = \| x - x_{\text{goal}} \|^2.
\end{equation}
The aim of static GRA is to find a control $u(t)$ (or parameters) such that $x(t) \to x^*$ and $\Phi \to 0$.

Equation of evolution (gradient descent):
\begin{equation}
    \dot{x} = -\nabla \Phi(x) + \text{noise (if present)}.
\end{equation}

\subsection{Applicability Criterion}

\textbf{Mathematical:} all eigenvalues of the linearization matrix $A = \frac{\partial f}{\partial x}\big|_{x^*}$ have negative real parts ($\Re(\lambda_i) < 0$). The system is asymptotically stable.

\textbf{Physical:} dissipation dominates in the system; entropy removal occurs faster than its production; the goal is to reach a single optimal state.

\subsection{Examples}
\begin{itemize}
    \item Thermostat (maintaining a given temperature).
    \item Robot position regulator.
    \item Sensor calibration.
    \item Traveling salesman problem (finding the shortest route).
\end{itemize}

\subsection{Conclusion}
Static GRA is applicable when the system can be completely calmed down without harming its functioning. It yields an exact, predictable solution.

\section{Cyclic GRA (Limit Cycle)}
\label{sec:cyclic}

\subsection{Mathematical Model}

The system tends to a limit cycle — a closed trajectory $\Gamma = \{x(t): x(t+T) = x(t)\}$ with minimal amplitude. The foam is defined as the root-mean-square deviation from some ``ideal'' cycle or as the oscillation amplitude:
\begin{equation}
    \Phi_{\text{cycle}} = \frac{1}{T} \int_0^T \| x(t) - x_{\text{ref}}(t) \|^2 \, dt.
\end{equation}
The goal is to minimize $\Phi_{\text{cycle}}$ while preserving self-oscillations (so that the cycle does not degenerate into a point).

Dynamics with control:
\begin{equation}
    \dot{x} = f(x) + u(t), \quad u(t) = -\nabla_x \Phi_{\text{cycle}} + \beta (\varphi - \varphi_0),
\end{equation}
where $\varphi$ is the cycle phase and $\varphi_0$ is the desired frequency.

\subsection{Applicability Criterion}

\textbf{Mathematical:} the linearized system at the stationary point has a pair of purely imaginary eigenvalues $\pm i\omega$, the rest have $\Re < 0$, and the first Lyapunov coefficient $l_1 < 0$ (Andronov–Hopf bifurcation). In this case a stable limit cycle is born from the point.

\textbf{Physical:} the system has an internal energy source (pumping) and dissipation, but they are balanced so that a self-oscillatory regime emerges. Complete stop is impossible or undesirable.

\subsection{Examples}
\begin{itemize}
    \item Heart rhythm (contraction cycle).
    \item Weight-driven clock.
    \item Belousov–Zhabotinsky chemical reactions.
    \item Breathing.
    \item Market price oscillations.
\end{itemize}

\subsection{Conclusion}
Cyclic GRA is necessary for systems that must ``breathe'' (oscillate) in order to function. It does not zero out the foam completely, but minimizes the amplitude and stabilizes the frequency.

\section{Chaotic GRA (Strange Attractor)}
\label{sec:chaotic}

\subsection{Mathematical Model}

The system exhibits deterministic chaos: trajectories are sensitive to initial conditions, aperiodic, but bounded in phase space (strange attractor). The foam is defined as the average energy of fluctuations or the Kolmogorov–Sinai entropy:
\begin{equation}
    \Phi_{\text{chaos}} = \langle \| x - \bar{x} \|^2 \rangle,
\end{equation}
where $\langle \cdot \rangle$ denotes time averaging (or ensemble averaging).

The aim of chaotic GRA is to minimize $\Phi_{\text{chaos}}$ by controlling parameters, or to transfer the system into a limit cycle or fixed point regime (chaos control).

Control equation (Pyragas method):
\begin{equation}
    \dot{x} = f(x) + K(x(t-\tau) - x(t)),
\end{equation}
where $K$ is a coefficient matrix and $\tau$ is the period of the desired cycle. By tuning $K$ one can stabilize an unstable cycle embedded in the chaos.

\subsection{Applicability Criterion}

\textbf{Mathematical:} the largest Lyapunov exponent $\lambda_1 > 0$ (the system is chaotic). If control can make $\lambda_1 \to 0$, cyclic or static GRA is applicable; otherwise, statistical chaotic GRA (minimization of averages) is used.

\textbf{Physical:} the system is strongly nonequilibrium, has many degrees of freedom, and is sensitive to initial conditions (weather, turbulence, market). Full prediction is impossible, but one can influence statistical characteristics.

\subsection{Examples}
\begin{itemize}
    \item Atmospheric convection (hurricanes).
    \item Fluid turbulence.
    \item Market crashes (volatility).
    \item Cardiac fibrillation (pathological chaos).
    \item Laser with chaotic modulation.
\end{itemize}

\subsection{Conclusion}
Chaotic GRA does not aim to eliminate foam completely, but only to reduce its statistical characteristics (variance, entropy). In some cases chaos control can transfer the system to a cycle or point regime.

\section{Hybrid GRA (Multilevel)}
\label{sec:hybrid}

\subsection{Mathematical Model}

Different GRA types are used at different hierarchy levels $l$. For example:
\begin{itemize}
    \item Lower levels (fast subsystems) – static GRA.
    \item Middle levels – cyclic GRA.
    \item Upper levels – chaotic GRA (or statistical).
\end{itemize}

Overall functional:
\begin{equation}
    J_{\text{hybrid}} = \sum_{l=0}^{K} w_l \cdot \Phi_{\text{type}(l)}^{(l)},
\end{equation}
where $\Phi_{\text{type}(l)}^{(l)}$ is the foam computed according to the rules corresponding to the chosen type for level $l$.

\subsection{Applicability Criterion}

Time hierarchy: fast variables can be stabilized, slow ones oscillate, and the slowest are chaotic. Such separation is common in complex systems (neural networks, economics, ecosystems).

\subsection{Examples}
\begin{itemize}
    \item AGI/ASI: fast signal processing (static), planning (cyclic), creativity (chaotic).
    \item Robot: posture stabilization (static), walking (cyclic), adaptation to new environments (chaotic search).
    \item Ecosystem: population cycle (cyclic), climatic fluctuations (chaotic), reserve management (static).
\end{itemize}

\subsection{Conclusion}
Hybrid GRA is the most general approach, combining the advantages of all three types at different scales. It requires a detailed analysis of time and parameter hierarchies.

\section{Comparative Table of GRA Types}

\begin{table}[h!]
\centering
\caption{Comparison of fundamental GRA types}
\label{tab:comparison}
\begin{tabular}{@{}lccc>{\raggedright}p{3cm}@{}}
\toprule
GRA Type & Attractor & Criterion & Foam $\Phi$ & Example \\
\midrule
Static & Fixed point & $\forall i:\ \Re(\lambda_i)<0$ & $\to 0$ & Thermostat \\
Cyclic & Limit cycle & $\lambda_{1,2}=\pm i\omega,\ l_1<0$ & Minimal amplitude & Heart \\
Chaotic & Strange attractor & $\lambda_1>0$ & Statistical variance & Hurricane \\
Hybrid & Mixed & Time hierarchy & Weighted sum & AGI/ASI \\
\bottomrule
\end{tabular}
\end{table}

\section{General Conclusions on GRA Types}

\begin{itemize}
    \item The choice of GRA type is determined by the dynamic properties of the system, not by the developer's whim. One cannot apply static GRA to a system that by its nature must oscillate (heart) or is chaotic (hurricane).
    \item Static GRA is the simplest and most studied, but its applicability is limited to systems where complete damping is desirable and possible.
    \item Cyclic GRA is essential for living systems, clocks, economic cycles, and any self-oscillatory processes.
    \item Chaotic GRA is the most complex case; it often reduces to statistical minimization or chaos control.
    \item Hybrid GRA is a promising path for creating AGI/ASI, where different levels of abstraction require different regimes.
\end{itemize}

\textbf{Formula for choosing the GRA type:}
\begin{equation}
    \text{GRA Type} = F\bigl( \{\Re(\lambda_i)\},\ l_1,\ \lambda_{\text{Lyapunov}},\ \text{time hierarchy} \bigr).
\end{equation}
\textbf{Practical advice:} always start with linearization and eigenvalue computation. If all $\Re(\lambda_i)<0$ – static GRA. If there is a pair of imaginary eigenvalues – check $l_1$ for cyclic. If there is a positive Lyapunov exponent – chaotic or hybrid is needed. Ignoring this rule leads to non-working algorithms.

\newpage

% =================== SET-THEORETIC EXTENSION ===================
\section{Set-Theoretic Representation of GRA-Zeroing Types}
\label{sec:set}

\subsection{The Set of GRA Types}

Define the set $\G$ of all fundamental GRA-zeroing types (excluding hybrid combinations):
\begin{equation}
    \G = \{ S, C, H \},
\end{equation}
where $S$ – static GRA (attractor – fixed point), $C$ – cyclic GRA (limit cycle), $H$ – chaotic GRA (strange attractor).

Hybrid types are represented as elements of the power set $\Pset(\G)$ or the Cartesian product $\G^n$ for hierarchical systems.

\subsection{Operations on Elements of the Set $\G$}

The following operations can be defined on $\G$, interpreted in the context of dynamical systems and foam control.

\subsubsection{Union $A \cup B$}

\textbf{Formally:} a binary operation on $\G$ with values in $\Pset(\G)$.

\textbf{Interpretation:} the system may operate in a regime where both GRA types are simultaneously applicable (or required) for different subsystems or phase-space domains. The union corresponds to an additive combination of foam functionals:
\begin{equation}
    \Phi_{A \cup B}(x) = \alpha \Phi_A(x) + \beta \Phi_B(x), \quad \alpha,\beta \geq 0,\ \alpha+\beta=1.
\end{equation}

\textbf{Example:} in a robot, the lower level (angle stabilization) uses $S$, the upper level (route planning with periodic traversal) uses $C$. The overall strategy is then $S \cup C$ with weights determined by the time hierarchy.

\textbf{Control selection formula:}
\begin{equation}
    u_{A \cup B}(t) = \gamma_A(t) u_A(t) + \gamma_B(t) u_B(t),
\end{equation}
where $u_A, u_B$ are controls synthesized according to the rules of types $A$ and $B$, and $\gamma_A, \gamma_B$ are activation functions depending on the system state.

\subsubsection{Intersection $A \cap B$}

\textbf{Interpretation:} the system possesses properties such that both types are simultaneously applicable (e.g., bistability or the presence of multiple attractors). The intersection defines a region in parameter space where the criteria of both types are jointly satisfied.

\textbf{Mathematically:} let $\mathcal{S}_A$ be the set of systems for which type $A$ is applicable (e.g., based on eigenvalue criteria). Then
\begin{equation}
    \mathcal{S}_{A \cap B} = \mathcal{S}_A \cap \mathcal{S}_B.
\end{equation}

Foam in such a system is minimized according to the intersection of functionals:
\begin{equation}
    \Phi_{A \cap B} = \max(\Phi_A, \Phi_B) \quad \text{or} \quad \sqrt{\Phi_A^2 + \Phi_B^2}.
\end{equation}

\textbf{Practical application:} if a system is near a bifurcation boundary (e.g., birth of a cycle from a fixed point), the control must simultaneously suppress noise and maintain oscillation amplitude within a safe range. The intersection $S \cap C$ provides a compromise solution.

\subsubsection{Difference $A \setminus B$}

\textbf{Interpretation:} application of type $A$ with an explicit prohibition against behavior characteristic of type $B$. For instance, $C \setminus S$ denotes cyclic GRA without degeneration into a fixed point (ensuring a minimum amplitude).

\textbf{Constraint formula:} a penalty term is introduced into the foam functional for approaching the undesired regime:
\begin{equation}
    \Phi_{A \setminus B} = \Phi_A + \lambda \cdot \dist(x, \mathcal{A}_B),
\end{equation}
where $\mathcal{A}_B$ is the attractor of type $B$, $\dist$ is the distance in phase space, and $\lambda \gg 1$.

\textbf{Example:} for cardiac rhythm, complete cessation ($S$) is forbidden; therefore $C \setminus S$ is used with a barrier function preventing the amplitude from falling below a threshold.

\subsubsection{Symmetric Difference $A \Delta B$}

\textbf{Interpretation:} selection between two mutually exclusive types depending on the system's operating regime. This forms the basis for switched hybrid systems.

\textbf{Formally:} control is defined piecewise:
\begin{equation}
    u(t) = \begin{cases}
        u_A(t), & \text{if } x \in \mathcal{D}_A, \\
        u_B(t), & \text{if } x \in \mathcal{D}_B,
    \end{cases}
\end{equation}
where $\mathcal{D}_A \cap \mathcal{D}_B = \emptyset$ and $\mathcal{D}_A \cup \mathcal{D}_B = \mathcal{D}$ is the working domain.

\textbf{Example:} during normal walking, a robot applies $C$ (cyclic GRA); upon detecting an obstacle, it switches to $S$ (static stabilization for a halt). This is a symmetric difference with event-triggered switching.

\subsubsection{Complement $\overline{A}$ with Respect to a Universe}

\textbf{Universe} $\U$ – the set of all possible methods for minimizing foam in dynamical systems, including non-GRA methods (e.g., stochastic control, adaptive control without an explicit attractor).

\textbf{Interpretation:} $\overline{A}$ denotes all approaches that are not of GRA type $A$.

\textbf{Practical significance:} determining the boundaries of applicability for GRA methods and identifying cases where alternative approaches are necessary (e.g., systems with non-Gaussian noise or unknown models).

\subsection{Cartesian Product and Hybrid Types}

Hybrid GRA, discussed in Section~\ref{sec:hybrid}, is formally represented as an element of a Cartesian product:
\begin{equation}
    \G_{\text{hybrid}}^{(L)} = \G^L = \underbrace{\G \times \G \times \dots \times \G}_{L \text{ times}},
\end{equation}
where $L$ is the number of hierarchical levels (or timescales). An element $(g_0, g_1, \dots, g_{L-1}) \in \G^L$ specifies the GRA type at each level.

\textbf{Foam functional of a hybrid system:}
\begin{equation}
    \Phi_{\text{hybrid}}(x) = \sum_{l=0}^{L-1} w_l \cdot \Phi_{g_l}(x^{(l)}),
\end{equation}
where $x^{(l)}$ are the variables of level $l$, and $w_l$ are weighting coefficients reflecting the time hierarchy.

\textbf{Operation of level composition:} if a given level employs not a single type but a combination, then $\G$ is replaced by $\Pset(\G)$.

\textbf{Example:} AGI/ASI architecture: \\
Level 0 (sensors) – $S$, \\
Level 1 (reflexes) – $C$, \\
Level 2 (planning) – $C \cup H$, \\
Level 3 (creativity) – $H$. \\
The full type is then the tuple $(S, C, C \cup H, H) \in (\Pset(\G))^4$.

\section{GRA-Zeroing as an Operator}
\label{sec:operator}

\subsection{Operator Definition}

Consider the space $\M$ of all controlled dynamical systems:
\begin{equation}
    \dot{x} = f(x, u), \quad x \in \R^n,\ u \in \R^m.
\end{equation}

To each system $(f, \text{goal})$ we can associate a GRA type $g \in \G$ (or its hybrid extension) according to a rule based on spectral properties of the linearization and Lyapunov exponents (see Sections~\ref{sec:static}--\ref{sec:chaotic}). This defines the \textbf{typing operator}:
\begin{equation}
    \T : \M \to \Pset(\G) \cup \G^L.
\end{equation}

Conversely, given a type $g$, one can synthesize a \textbf{control operator}:
\begin{equation}
    \U_g : \M_{\text{open}} \to \M_{\text{closed}},
\end{equation}
which maps an open-loop system to a closed-loop system with a controller implementing GRA-zeroing of type $g$.

\textbf{Composition of operators:}
\begin{equation}
    \U_{g_1} \circ \U_{g_2}
\end{equation}
denotes the sequential application of two strategies (e.g., first stabilizing a point, then introducing small oscillations to create a cycle). This is possible if the corresponding subsystems are separated in time or space.

\textbf{Identity operator} $\I$ corresponds to the absence of control (the system retains its original dynamics). It does not belong to $\G$ but is part of the universe $\U$.

\textbf{Inverse operator} $\U_g^{-1}$ (if it exists) – transition from a closed-loop system back to the open-loop system, i.e., identification of the original dynamics from observed behavior under the controller.

\subsection{Algebra of GRA Operators}

On the set of operators $\{\U_g \mid g \in \G^*\}$, one can define operations induced by the set-theoretic operations on types:

\begin{itemize}
    \item \textbf{Sum of operators:} $\U_A \oplus \U_B = \U_{A \cup B}$ (additive control).
    \item \textbf{Product of operators:} $\U_A \otimes \U_B = \U_{A \cap B}$ (joint application with constraints).
    \item \textbf{Inversion:} $\U_{\overline{A}}$ – an operator implementing any method except GRA type $A$.
\end{itemize}

\textbf{Practical meaning:} This algebra allows formal description of complex controllers as combinations of basic GRA strategies. For example, a controller for a walking robot can be written as:
\begin{equation}
    \U_{\text{robot}} = (\U_S \oplus \U_C) \otimes \U_{C \setminus S},
\end{equation}
meaning: simultaneous stabilization and cycling, but with a prohibition against cycle arrest.

\section{Practical Applications of the Set-Theoretic Approach}
\label{sec:applications}

\subsection{Automatic System Classification}

Given a numerical model of a system, one can compute:
\begin{itemize}
    \item eigenvalues of the linearized matrix $\{\lambda_i\}$,
    \item first Lyapunov coefficient $l_1$,
    \item largest Lyapunov exponent $\lambda_{\max}$.
\end{itemize}

From these data, the mapping $\T(f)$ is constructed. The set $\G$ serves as a label space for a classifier. For example:

\begin{table}[h!]
\centering
\caption{Rules for classifying a system by GRA type}
\label{tab:classification}
\begin{tabular}{@{}lc@{}}
\toprule
Condition & GRA Type \\
\midrule
$\forall i:\ \Re(\lambda_i) < 0$ & $S$ \\
$\exists i:\ \Re(\lambda_i) = 0,\ l_1 < 0$ & $C$ \\
$\lambda_{\max} > 0$ & $H$ \\
\bottomrule
\end{tabular}
\end{table}

\subsection{Designing Hierarchical Controllers}

When building AGI/ASI architectures, each level of abstraction is assigned an element from $\Pset(\G)$ or $\G^L$. Union and intersection operations enable a formal description of level interactions. For instance, the consciousness level may use $H \cap C$ – chaotic search with periodic attentional focusing.

\subsection{Verification and Debugging}

Using difference and complement operations, one can formally specify \textbf{forbidden regimes} (e.g., $\overline{S}$ at the respiratory level – prohibition of respiratory arrest). This simplifies the specification of safety requirements.

\subsection{Adaptive Switching of Types}

Based on symmetric difference, hysteresis-based switching between $S$ and $C$ can be implemented depending on external conditions (e.g., economic policy: stabilization in crisis, cyclic regulation in growth).

\section{Conclusions}
\label{sec:conclusions}

\begin{enumerate}
    \item \textbf{Set-Theoretic Nature of GRA:} The fundamental GRA-zeroing types form a discrete set $\G$ on which Boolean operations are naturally defined. These operations have clear physical interpretations in terms of combining foam functionals and imposing constraints on attractors.
    \item \textbf{Hybridity as a Cartesian Product:} Hierarchical and multi-level systems are described by elements of $\G^L$ or $\Pset(\G)^L$, enabling the formal construction of complex control strategies.
    \item \textbf{Operator Algebra:} GRA-zeroing can be represented as an operator acting on the space of dynamical systems. Composition, sum, and product of operators reflect the combination of foam-minimization strategies.
    \item \textbf{Practical Value:} The set-theoretic approach provides:
    \begin{itemize}
        \item a unified language for specifying controller requirements;
        \item a formal basis for automated synthesis of hybrid controllers;
        \item a tool for analyzing the compatibility of different control methods.
    \end{itemize}
\end{enumerate}

\textbf{Final formula for the GRA-type selection operator with multiple constraints:}
\begin{equation}
    g^* = \arg\min_{g \in \mathcal{G}} \left\{ \Phi_g(x) + \sum_{i} \mu_i \cdot \mathbb{I}[x \notin \mathcal{D}_{g_i}] \right\},
\end{equation}
where $\mathcal{G} \subseteq \Pset(\G)$ is the admissible subset of types, $\mathbb{I}$ is the indicator function for forbidden regions, and $\mu_i$ are large penalty coefficients.

This approach transforms GRA-zeroing from a collection of disparate methods into a rigorous algebraic system, ready for application in the most complex control problems, including the development of strong artificial intelligence.

% Acknowledgments (optional)
\section*{Acknowledgments}
This work was supported by RFBR grant No.~XX-XX-XXXXX.

% Bibliography
\begin{thebibliography}{9}
\bibitem{pyragas} Pyragas K. Continuous control of chaos by self-controlling feedback // Physics Letters A, 1992.
\bibitem{kuznetsov} Kuznetsov Yu.A. Elements of Applied Bifurcation Theory. — Springer, 2004.
\bibitem{strogatz} Strogatz S.H. Nonlinear Dynamics and Chaos. — Westview Press, 2014.
\bibitem{andronov} Andronov A.A., Vitt A.A., Khaikin S.E. Theory of Oscillators. — Pergamon Press, 1966.
\end{thebibliography}

\end{document}